\documentclass{ximera}



\begin{document}
	\author{Wim Obbels}
	\xmtitle{Tweedegraadsvergelijkingen in de complexe getallen}{}
	\label{xim:cmplx_tweedegraadsvergelijkingen}

    Wiskundigen zijn dikwijls erg enthousiast over complexe getallen omdat ze de wiskunde vereenvoudigen.
    \\
    Dit lijkt voor sommigen misschien moeilijk te geloven, maar we zullen het proberen aan te tonen door het oplossen van tweedegraadsvergelijkingen opnieuw te bestuderen nu we stilaan vertrouwd raken met de complexe getallen.

    In de tweede graad  werd reeds gezien hoe je tweedegraadsvergelijkingen oplost: 

    \begin{expandable}{remark}{Vierkantsvergelijkingen $ax^2+bx+c=0$ oplossen in $\R$.}\nl

        \begin{proposition}\nl

            Een (reële) tweedegraadsvergelijking van de vorm $ax^2+bx+c$, met $a,b,c\in\R$ en $a\neq0$
            heeft als \textbf{discriminant} het (reële) getal  \important{D=b^2-4ac}, 
            en heeft als oplossingen

            \begin{tabular}{lll}
                als $D<0$: & geen reële oplossingen \\
                als $D=0$: & precies een reële oplossing, namelijk \important{x_1=-\dfrac{b}{2a}} \\
                als $D>0$: & precies twee reële oplossingen, namelijk 
                             \important{x_1=\dfrac{-b+\sqrt{D}}{2a}} en \important{x_2=\dfrac{-b-\sqrt{D}}{2a}}
            \end{tabular}

            Bovendien zijn de volgende uitspraken equivalent:

            \begin{tabular}{lll}
                (a) & $x_1$ en $x_2$ zijn oplossingen van de vergelijking $ax^2+bx+c=0$ \\
                (b) & $x_1$ en $x_2$ zijn nulpunten  van de functie $f(x)=ax^2+bx+c$ \\
                (c) & $x_1$ en $x_2$ zijn snijpunten van de kromme $y=ax^2+bx+c$ met de $x$-as \\
                (d) & $ax^2+bx+c = a(x-x_1)(x-x_2)$ & (ontbinden in factoren)\\
                (e) & $x_1+x_2 = -\dfrac{b}{a}$ en $x_1\cdot x_2 = \frac{c}{a}$ & (som en product van de wortels)
            \end{tabular}

        \end{proposition}

    \end{expandable}

    Een negatieve discriminant stelde in de tweede graad een probleem, want uit negatieve getallen kan je in de reële getallen geen vierkantswortels trekken.
    Bij complexe getallen verdwijnt dat probleem:

    \begin{proposition}\nl

        Een (reële) tweedegraadsvergelijking van de vorm $ax^2+bx+c$, met $a,b,c\in\R$ en $a\neq0$
        heeft altijd twee complex toegevoegde oplossingen, namelijk 
        \begin{center}
            \important{x_1=\dfrac{-b+\sqrt{b^2-4ac}}{2a}} en \important{x_2=\dfrac{-b-\sqrt{b^2-4ac}}{2a}}
        \end{center}
        die echter samenvallen als $b^2-4ac=0$, en dan gelijk worden aan $x_1=x_2=\dfrac{-b}{2a}$.
    \end{proposition}
    Het begrip discriminant wordt dus in zekere zin overbodig (of minstens veel minder belangrijk).


    \begin{remark}{Een negatief getal onder de wortel schrijven?}

        Door bovenstaande formule in te vullen kan je een negatief getal onder de wortel uitkomen. 
        \textbf{Een negatief getal onder de wortel is niet goed gedefinieërd:} de rekenregel \(\sqrt{\square \square} = \sqrt{\square}\sqrt{\square}\) zorgt meteen voor een tegenstrijdigheid: 
    %underbrackets met -1 invoegen 

    Zo is $(\sqrt{-1})^2 = -1 $ als we de wortel en het kwadraat schrappen, maar als we de rekenregel toepassen vinden we 
    $$\sqrt{-1}^2 = \sqrt{-1} \cdot \sqrt{-1} = \sqrt{-1 \cdot -1} = \sqrt{1} = 1 $$

    Omwille van deze tegenstrijdigheid maken we volgende afspraak: 
    \begin{center} 
        \textbf{We schrijven \(\sqrt{\square}\) enkel voor \( \square \in \R \)}. 
    \end{center}

    In de oplossingen vullen we de formule uit bovenstaande eigenschap in en zal je soms in een tussenstap een negatief getal onder de wortel zien: deze uitdrukking beschouwen we louter als een symbolische notatie in een tussenstap.   
    \textbf{Laat een dergelijke uitdrukking nooit staan in je eindoplossing aangezien het geen goed gedefinieërde wiskundige betekenis heeft.}        
    \end{remark}

    \begin{example} 
        \begin{question} De wortels van de vergelijking $x^2+4=0$ zijn $x_1=2i$ en $x_2=-2i$.\end{question}
        \begin{question} De wortels van de vergelijking $x^2+x+1=0$ zijn $x_{1,2}=\dfrac{-1\pm i\sqrt{3}}{2}$\end{question}
    \end{example}

    \begin{exercise} Bereken de wortels van volgende vergelijkingen.
        \begin{question} $x^2+2x+3=0$
        \begin{oplossing}\nl
            Hier is $a=1$, $b=2$ en $c=3$, dus de discriminant is $D=b^2-4ac=4-12=-8$.
            De wortelformule geeft
            $$
            x_{1,2}=\dfrac{-2\pm\sqrt{-8}}{2}=\dfrac{-2\pm 2i\sqrt{2}}{2}=-1\pm i\sqrt{2}\,.
            $$
        \end{oplossing}
        \end{question}
        \begin{question} $2x^2+3x+4=0$
        \begin{oplossing}\nl
            Hier is $a=2$, $b=3$ en $c=4$, dus $D=9-32=-23$.
            Dus
            $$
            x_{1,2}=\dfrac{-3\pm\sqrt{-23}}{4}=\dfrac{-3\pm i\sqrt{23}}{4}\,.
            $$
        \end{oplossing}
        \end{question}
    \end{exercise}

    \begin{exercise} Ontbind volgende veeltermen in lineaire factoren.
        \begin{question} $x^2+2x+3$
        \begin{oplossing}\nl
            Uit de vorige oefening volgt dat de nulpunten van $x^2+2x+3$ gelijk zijn aan $-1\pm i\sqrt{2}$.
            Daarom is
            $$
            x^2+2x+3=(x-(-1+i\sqrt{2}))(x-(-1-i\sqrt{2}))
            =(x+1-i\sqrt{2})(x+1+i\sqrt{2})\,.
            $$
        \end{oplossing}
        \end{question}
        \begin{question} $2x^2+3x+4$
        \begin{oplossing}\nl
            De nulpunten van $2x^2+3x+4$ zijn $\dfrac{-3\pm i\sqrt{23}}{4}$ (zie vorige oefening).
            Omdat de coëfficiënt van $x^2$ gelijk is aan $2$, geldt
            $$
            2x^2+3x+4
            =2\left(x-\dfrac{-3+i\sqrt{23}}{4}\right)\left(x-\dfrac{-3-i\sqrt{23}}{4}\right)\,.
            $$
        \end{oplossing}
        \end{question}
    \end{exercise}

\end{document}